\chapter{Zakończenie}

Rosnąca złożoność środowisk ICT w organizacjach nadzorowanych, wzmocniona przez wejście w życie rozporządzenia DORA i implementację dyrektywy NIS2, stworzyła wyraźną lukę pomiędzy wymaganiami regulacyjnymi a dostępnymi narzędziami ich weryfikacji. Istniejące metodyki audytu (ISO 27001, COBIT, NIST CSF) dostarczają ogólnych ram kontrolnych, lecz nie oferują skwantyfikowanego, powtarzalnego mechanizmu oceny ryzyka uwzględniającego jednocześnie strukturę łańcucha dostaw ICT, klasyfikację przetwarzanych danych, specyfikę jurysdykcji dostawców chmurowych oraz wagę poszczególnych wymagań regulacyjnych. Niniejsza praca odpowiada na ten deficyt, proponując metodologię ORION (Ocena Ryzyka Informatycznego Organizacji Nadzorowanych), zaprojektowaną z myślą o organizacjach objętych regulacjami sektorowymi i unijną legislacją z zakresu cyfrowej odporności operacyjnej.

Centralnym elementem ORION jest model grafu ważonego $G = (V, E, W)$, w którym wierzchołki reprezentują obiekty audytowe trzech typów (PODMIOT, SYSTEM INFORMATYCZNY, DOSTAWCA), a krawędzie kierowane odzwierciedlają zależności usługowe z wagami odpowiadającymi krytyczności obsługiwanych systemów. Model ten umożliwia propagację ryzyka wzdłuż łańcucha dostaw i automatyczne rozszerzanie zakresu audytu o poddostawców w przypadku aktywacji triggera TIA (Third-party Impact Analysis), co bezpośrednio adresuje wymogi art. 28 DORA dotyczące mapowania zależności zewnętrznych. Mechanizm walidacji krzyżowej (każde pytanie zadawane niezależnie trzem perspektywom: compliance/prawnej, IT oraz użytkownika) ogranicza ryzyko konfabulacji i formalnego spełniania wymagań bez ich realnego wdrożenia. Wynik końcowy $S_{total}$ jest wielkością skwantyfikowaną, wyznaczaną na podstawie jednolitego wzoru ważonego z jasno zdefiniowaną skalą pięciostopniową, co czyni ocenę porównywalną między audytami i organizacjami.

Przeprowadzony audyt pilotażowy systemu CTN w instytucji bankowej potwierdził wykonalność i spójność metodologii w warunkach rzeczywistych. Spośród 62 pytań kwalifikacyjnych 58 uzyskało odpowiedź pozytywną. Cztery niezgodności (dwie obowiązkowe dotyczące rotacji haseł i braku testu exit planu, dwie zalecane z zakresu bezpieczeństwa fizycznego) przełożyły się na wynik $S_{total} = 0{,}951$, klasyfikowany jako ZGODNY w przedziale $[0{,}85;\ 1{,}0]$. Żadna sekcja nie aktywowała flagi blokującej. Trigger TIA nie został wyzwolony, ponieważ dostawca systemu działa w jurysdykcji EOG, a wszystkie umowy spełniały wymagania DORA w zakresie klauzul audytowych i prawa do wypowiedzenia. Wynik wskazuje, że organizacja prowadzi zarządzanie ryzykiem ICT na dojrzałym poziomie, jednocześnie precyzyjnie wskazując dwa obszary wymagające korekty w ramach działań naprawczych.

Jedną z celowych właściwości projektowych ORION jest jego adaptowalność. Choć poniższy przypadek użycia dotyczy sektora finansowego (banku nadzorowanego przez KNF), architektura metodologii nie zakłada żadnych elementów specyficznych wyłącznie dla bankowości. Biblioteka sekcji (bezpieczeństwo techniczne, chmura, dane poufne, ciągłość działania, zarządzanie dostawcami) stanowi zbiór bloków, które mogą być selektywnie włączane lub wyłączane na podstawie atrybutów obiektu audytowego i katalogu regulacji mającego zastosowanie do danej branży. Szpital przetwarzający dane medyczne klasy D3 aktywuje sekcje dotyczące ochrony danych poufnych z inną wagą niż operator logistyczny, lecz korzysta z tego samego aparatu formalnego i tych samych procedur zbierania dowodów. Dostawca energii objęty NIS2 jako podmiot kluczowy, podmiot administracji publicznej prowadzący rejestr ludności, platforma przetwarzania danych medycznych, czy firma fintech objęta DORA, wszystkie te organizacje mają odmienną specyfikę ryzyka, lecz wspólny mianownik: konieczność wykazania organom nadzoru, że ryzyko ICT jest zarządzane w sposób mierzalny i udokumentowany. ORION dostarcza do tego celu jednolity język.

Żaden system oceny ryzyka nie może być lepszy niż jakość danych, na których operuje, a ta zależy przede wszystkim od kompetencji i rzetelności audytora. Kluczowe znaczenie ma konsekwentne przestrzeganie zasady niezależności perspektyw w walidacji krzyżowej: zbyt swobodna interpretacja odpowiedzi jednej perspektywy jako potwierdzenia dla pozostałych wypacza wynik i podważa sens całego mechanizmu. Równie istotna jest wytrwałość w dążeniu do weryfikacji dowodów (logi, konfiguracje, umowy, wyniki testów) zamiast poprzestania na deklaracjach słownych. Przypadek wykrycia braku MFA w trakcie audytu CTN, niezgodności naprawionej jeszcze przed zamknięciem kwestionariusza, ilustruje, jak rzetelna inspekcja może mieć bezpośredni, pozytywny wpływ na bezpieczeństwo systemu. Metodologia ORION tworzy ramy do przeprowadzenia takiej inspekcji, lecz jej realna wartość zależy od gotowości audytora do konfrontowania się z trudnymi obserwacjami i formułowania wniosków bez ulegania presji organizacyjnej.

Autor identyfikuje kilka kierunków dalszego rozwoju metodologii. Po pierwsze, automatyzacja zbierania i walidacji dowodów technicznych (logi SIEM, wyniki skanerów podatności, raporty certyfikatów TLS) może istotnie skrócić czas trwania fazy kwestionariuszowej i ograniczyć ryzyko błędów ludzkich. Po drugie, mechanizm propagacji ryzyka w grafie ICT może zostać rozszerzony o modele probabilistyczne uwzględniające historyczne dane o incydentach, co pozwoliłoby na szacowanie oczekiwanej straty operacyjnej powiązanej z konkretnym poddostawcą. Po trzecie, opracowanie branżowych szablonów pytań (dla sektora zdrowotnego, energetycznego, administracji publicznej) pozwoliłoby skrócić etap konfiguracji audytu przy zachowaniu wspólnego aparatu obliczeniowego. Po czwarte, integracja z rejestrami CMDB i systemami GRC (Governance, Risk and Compliance) umożliwiłaby ciągłą, automatyczną aktualizację grafu ICT bez konieczności inicjowania pełnego audytu przy każdej zmianie w środowisku.

Metodologia ORION nie pretenduje do bycia jedynym ani wyczerpującym narzędziem oceny ryzyka ICT w organizacjach nadzorowanych. Jest natomiast propozycją wypełnienia konkretnej luki: braku ustrukturyzowanej, skwantyfikowanej i regulacyjnie zakorzenionej metody łączącej audyt organizacyjny z analizą łańcucha dostaw w jednolitym modelu grafowym. Zaprezentowany przypadek użycia potwierdza, że taki model jest wykonalny w praktyce i generuje wnioski działające dwukierunkowo: dając organom nadzoru czytelny wskaźnik zgodności, a organizacji audytowanej precyzyjny priorytet działań naprawczych. Dalszy rozwój metodologii, jej walidacja na większej próbie organizacji i ewentualna standaryzacja w postaci normy branżowej lub wytycznych regulatora, pozostają zadaniami otwartymi i stanowią naturalny kierunek kontynuacji niniejszych badań.
